% entrevista_abuelo.tex
% Cuestionario impreso para que José Luis Guerrero Bobadilla lo conteste a mano.
% Blanco y negro, Times (TeX Gyre Termes) 12 pt.  Compilar:  latexmk -lualatex entrevista_abuelo.tex
\documentclass[12pt]{article}
\usepackage[spanish,es-noquoting]{babel}
\usepackage{fontspec}
\setmainfont{texgyretermes}[Extension=.otf, UprightFont=*-regular, BoldFont=*-bold,
  ItalicFont=*-italic, BoldItalicFont=*-bolditalic]
\usepackage[letterpaper, margin=2.2cm, bottom=2.4cm]{geometry}
\usepackage{enumitem}
\usepackage{needspace}
\usepackage{fancyhdr}
\usepackage{lastpage}

\pagestyle{fancy}\fancyhf{}
\renewcommand{\headrulewidth}{0pt}
\fancyfoot[C]{\small Página \thepage\ de \pageref{LastPage}}
\setlength{\parindent}{0pt}
\setlength{\parskip}{6pt}

% Renglones para escribir a mano (separación 0.95 cm)
\newcount\cuentarenglon
\newcommand{\renglones}[1]{\par\cuentarenglon=0
  \loop\ifnum\cuentarenglon<#1 \advance\cuentarenglon by 1
    \vspace{0.85cm}\hrule height 0.4pt
  \repeat\par}

% Pregunta numerada con su espacio de respuesta; va completa en una página
\newcounter{pregunta}
\newcommand{\pregunta}[2][6]{\stepcounter{pregunta}%
  \par\medskip\noindent\begin{minipage}{\linewidth}\setlength{\parskip}{0pt}%
  \textbf{\thepregunta.}~#2\par\renglones{#1}\end{minipage}\par}
\newcommand{\bloque}[1]{\par\needspace{7cm}\bigskip{\large\bfseries #1}\par\smallskip\hrule\par\nobreak}

\begin{document}

\begin{center}
  {\Large\bfseries Cuestionario: las maquiladoras en la frontera de Sonora}\\[4pt]
  {\large Testimonio de José Luis Guerrero Bobadilla}\\[2pt]
  Presidente municipal de Naco, Sonora (1976--1979)
\end{center}
\medskip

Querido abuelo:

Estoy escribiendo mi ensayo final para la materia de Contexto Regional de la Universidad de Sonora. El tema es la llegada de las \textbf{maquiladoras} a la frontera de Sonora (Nogales, Agua Prieta y Naco) entre los años sesenta y setenta. Como usted vivió esa época y además fue presidente municipal de Naco, sus recuerdos son una fuente muy valiosa.

\textbf{Cómo contestar:}
\begin{itemize}[nosep, leftmargin=1.5em]
  \item Escriba con sus propias palabras; no hay respuestas buenas ni malas.
  \item Si no recuerda algo, puede escribir «no recuerdo» o dejarlo en blanco.
  \item Si le falta espacio, siga escribiendo en las hojas en blanco del final e indique el número de la pregunta.
  \item Nombres de personas, empresas, lugares, fechas y cantidades (aunque sean aproximadas) ayudan mucho.
\end{itemize}

\textbf{Cómo se usarán sus respuestas:} algunas frases suyas aparecerán como citas textuales en el ensayo, con su nombre, como testimonio de alguien que vivió los hechos.

\medskip
Acepto que mis respuestas se citen con mi nombre en el ensayo.\par\medskip
Firma: \rule{6.5cm}{0.4pt}\hfill Fecha: \rule{4.5cm}{0.4pt}

% ------------------------------------------------------------------
\bloque{A. Datos generales}
\pregunta[2]{¿En qué año y en qué lugar nació?}
\pregunta[2]{¿Desde qué año vive en Naco? Si llegó de otro lugar, ¿por qué se vino a Naco?}
\pregunta[4]{Antes de ser presidente municipal, ¿a qué se dedicaba? ¿Dónde trabajó?}

% ------------------------------------------------------------------
\bloque{B. Naco antes de las maquiladoras (años sesenta)}
\pregunta{¿De qué vivía la gente de Naco en los años sesenta? (Por ejemplo: la mina de Cananea, la ganadería, el comercio con Naco, Arizona, la aduana, el campo.)}
\pregunta{¿Había trabajo suficiente en Naco? ¿Mucha gente cruzaba a trabajar a Estados Unidos o se iba a otras ciudades?}
\pregunta[5]{En 1964 terminó el Programa Bracero y muchos trabajadores regresaron de Estados Unidos. ¿Recuerda si eso se sintió en Naco? ¿Qué pasó con ellos?}

% ------------------------------------------------------------------
\bloque{C. La llegada de las maquiladoras (1965--1975)}
\pregunta[5]{¿Cuándo escuchó hablar por primera vez de las maquiladoras? ¿Qué se decía de ellas en el pueblo?}
\pregunta[8]{¿Hubo maquiladoras en Naco? Si las hubo: ¿recuerda sus nombres, qué fabricaban, dónde estaban y cuánta gente trabajaba en ellas?}
\pregunta[5]{¿La gente de Naco iba a trabajar a las maquiladoras de Agua Prieta, Cananea o Nogales? ¿Cómo se transportaban?}
\pregunta{¿Quiénes trabajaban en las maquiladoras: más hombres o más mujeres, jóvenes o adultos? ¿Cómo eran los sueldos y las condiciones de trabajo?}
\pregunta{¿Cambió el pueblo con la llegada de las maquiladoras? (Comercio, población, vivienda, servicios como agua, luz y calles.)}

% ------------------------------------------------------------------
\bloque{D. La crisis de 1974--1975}
\textit{Según la investigadora Cirila Quintero, durante la crisis de 1974--1975 cerraron 34 empresas entre Nogales, Agua Prieta y Naco por la recesión en Estados Unidos.}
\pregunta{¿Recuerda esos cierres? ¿Qué pasó con los trabajadores que se quedaron sin empleo?}
\pregunta{¿A los trabajadores les pagaron lo que marcaba la ley (indemnización)? ¿Hubo sindicatos, protestas o huelgas?}

% ------------------------------------------------------------------
\bloque{E. Su periodo como presidente municipal (1976--1979)}
\pregunta{¿Cómo llegó a ser presidente municipal? ¿Cuáles eran los problemas más importantes de Naco cuando empezó su administración?}
\pregunta[8]{¿Atraer empresas o maquiladoras a Naco era una prioridad para usted? ¿Qué se hizo para lograrlo? (Por ejemplo: ofrecer terrenos, quitar o reducir impuestos, poner servicios, construir naves industriales.)}
\pregunta{¿Cómo era la relación con el gobierno del estado (gobernador Alejandro Carrillo Marcor) y con el gobierno federal? ¿Le dieron apoyo para traer industria o empleo?}
\pregunta{¿Tuvo trato directo con empresarios de Estados Unidos o de México que quisieran instalarse en Naco? ¿Qué pedían a cambio?}
\pregunta{En 1976 el peso se devaluó (el dólar pasó de 12.50 a cerca de 20 pesos). ¿Cómo afectó eso a Naco, a sus comerciantes y a sus trabajadores?}
\pregunta{¿Qué ventajas o desventajas tenía Naco frente a Agua Prieta o Nogales para atraer industria?}
\pregunta{¿Qué logro de su administración relacionado con el empleo o la economía del pueblo recuerda con más orgullo?}
\pregunta{¿Qué le hubiera gustado lograr y no se pudo? ¿Por qué?}

% ------------------------------------------------------------------
\bloque{F. Mirando hacia atrás}
\pregunta[8]{Viéndolo hoy, ¿las maquiladoras fueron buenas o malas para Naco y para la frontera? ¿Por qué?}
\pregunta{¿Quién cree que ganó más con las maquiladoras: los trabajadores, los empresarios o el gobierno?}
\pregunta{Si pudiera volver a esos años como presidente municipal, ¿qué haría diferente?}
\pregunta[8]{¿Tiene alguna anécdota o recuerdo de esa época (de las maquiladoras, de la crisis o de su administración) que quiera contarme?}
\pregunta[6]{¿Hay algo más que quiera agregar?}

% ------------------------------------------------------------------
\newpage
\bloque{Hojas adicionales}
Si le faltó espacio, continúe aquí. Escriba el número de la pregunta antes de su respuesta.
\renglones{22}
\newpage
\renglones{24}

\vfill
\begin{center}\itshape ¡Muchas gracias, abuelo!\end{center}
\end{document}
